\section{Related Work}
\label{sec:related}

\para{Attractor dynamics in neural networks.}
The study of attractors in trained neural networks draws on dynamical systems theory \citep{strogatz2015nonlinear}. \citet{beer2018interpreting} introduced methods for extracting attractor dynamics from recurrent networks by identifying invariant sets in phase space. More recently, \citet{chytas2025concept} formalized \emph{concept attractors} in large language models, showing that Transformer layers implement iterated function systems that map semantically related prompts to concept-specific fixed points in latent space. They demonstrated that different concept families form attractors at different layers---for example, programming concepts converge at layer 19 while natural language concepts converge at layer 27 in Llama 3.1 8B. \citet{wang2025unveiling} extended this line of work by showing that LLMs converge to stable 2-period attractor cycles under successive paraphrasing, a pattern that persists across languages, text lengths, and sampling temperatures. Both studies, however, examine only Transformer architectures, leaving open whether attractor formation is architecture-dependent. Our work directly addresses this gap.

\para{Architectural differences between RNNs and Transformers.}
\citet{fukushima2023comparing} compared RNN and Transformer generalization on 2D dynamical systems, finding that RNNs converge to correct attractor dynamics with only 2--4 training sequences while Transformers require 25--50 sequences---a difference attributed to the recurrent inductive bias of shared weights. \citet{wen2024rnns} proved a fundamental representation gap: RNNs with chain-of-thought have $O(\log n)$ bit memory, exponentially weaker than Transformers, due to a bottleneck in in-context retrieval. Adding a single Transformer layer to an RNN closes this gap, suggesting a sharp architectural boundary. \citet{barak2024transformers} bridged the two families by showing that Transformer attention can be viewed as a multi-state RNN, providing a theoretical framework for understanding when behaviors might converge. Unlike these works, which focus on capability differences, we study how architectural inductive biases shape the \emph{default outputs} that emerge from training on identical data.

\para{Mode collapse and persona convergence.}
The tendency of aligned language models to converge on similar default responses has been documented in several concurrent works. \citet{sycophancy2025} showed that LLMs develop systematic sycophantic behavior where user opinions actively suppress learned knowledge in later layers---a form of attractor dynamics in the behavioral space. \citet{mode_collapse2025} studied mode collapse under verbalized sampling, finding that RLHF constrains models from generalizing over authorship and produces narrower output distributions. These studies attribute convergence primarily to training data overlap and alignment procedures. Our work isolates the architectural component, showing that even without alignment, different architectures trained on identical data develop different default behaviors. This complements the existing explanations by identifying a third, previously unstudied driver of persona collapse.
